\documentclass[11pt,oneside,openany]{memoir}
\usepackage{fontspec}
\usepackage{polyglossia}
\setmainlanguage{french}
\setmainfont{TeX Gyre Pagella}
\setsansfont{TeX Gyre Heros}
\setmonofont{Latin Modern Mono}
\usepackage[final]{microtype}
\usepackage{csquotes}
\usepackage{hyperref}
\usepackage{verse}
\setstocksize{230mm}{155mm}
\settrimmedsize{\stockheight}{\stockwidth}{*}
\setlrmarginsandblock{23mm}{22mm}{*}
\setulmarginsandblock{24mm}{22mm}{*}
\checkandfixthelayout
\setlength{\parindent}{1.25em}
\setlength{\parskip}{0pt}
\renewcommand{\footnotesize}{\small}
\begin{document}
\title{Histoire éditoriale et humanités}
\author{Claire Martin}
\maketitle

\chapter*{Liminaires}
\addcontentsline{toc}{chapter}{Liminaires}
Ce volume propose un cadre de lecture pour la collection.

\chapter{Chapitre 1. Méthodes}
Les archives de la maison éditoriale comportent des mentions en\textsc{petites capitales}et des datations du XIX\textsuperscript{e}siècle.

L'analyse impose une note critique\footnote{Voir l'édition\textit{princeps}et ses variantes.}pour situer le témoin.

\begin{quote}\small
Cette citation longue en prose sert de bloc de démonstration pour la future épreuve de composition.
\end{quote}

\begin{verse}
Le texte se plie au rythme\\
Et la page retient le vers\\
\end{verse}

\chapter*{Conclusion}
\addcontentsline{toc}{chapter}{Conclusion}
La chaîne TEI vers LaTeX doit rester explicable et traçable.

\chapter*{Bibliographie}
\addcontentsline{toc}{chapter}{Bibliographie}
\section*{Bibliographie}
\addcontentsline{toc}{section}{Bibliographie}
\noindent Martin, Claire,\textit{Composer le savoir}, Rouen, 2022.\par
\noindent Durand, Paul,\textit{Éditer les sources}, Paris, 2019.\par

\end{document}
